% META: {"id":"1SPE-SUITES-CO-032","chapitre":"1SPE-SUITES","type_objet":"corrige","exercice_ref":"1SPE-SUITES-EX-032","capacites":["1SPE-SUITES-C3","1SPE-SUITES-C5","1SPE-SUITES-C6"],"capacites_codes":["C3","C5","C6"],"sources_inspiration":[],"status":"approved","origin":"generated","status_history":["generated","audited","accepted","approved"],"release_acceptance":"RELEASE_OWNER_BATCH_ACCEPTANCE","acceptance_closure_digest":"sha256:9b3ccf9a81c5520fbb7e03b7b2d3e7bf2057a02834b6d908bf3be5f49f3b553f","acceptance_receipt_digest":"sha256:4fad86f0282e0fa4e0419cca1ad6379ae7af5a280c04a4a90880f90a1c044242"}
% BEGIN-VERIFY
% from sympy import Rational, simplify, symbols
% n = symbols('n', integer=True, nonnegative=True)
% C0 = 800
% q = Rational(103, 100)
% C = C0 * q**n
% assert C.subs(n, 0) == 800
% ratio = simplify(C.subs(n, n+1) / C)
% assert ratio == Rational(103, 100)
% assert ratio > 1
% # C_5
% C5 = float(C0 * (1.03)**5)
% assert abs(C5 - 927.42) < 0.01
% # Seuil strict C_n > 1200 : 800 * 1.03^n > 1200 => 1.03^n > 1.5
% # 1.03^13 ~ 1.469 < 1.5, 1.03^14 ~ 1.513 > 1.5 => n = 14
% assert 800 * 1.469 < 1200
% assert 800 * 1.513 > 1200
% END-VERIFY
\begin{corrige}{1SPE-SUITES-EX-032}

\textbf{Question 1 — Relation de récurrence et nature géométrique.}

Chaque année, le capital est multiplié par $(1 + 0{,}03) = 1{,}03$. Donc :
\[
  C_{n+1} = 1{,}03 \times C_n, \quad C_0 = 800.
\]

D'après la relation de récurrence $C_{n+1}=1{,}03C_n$, chaque terme
s'obtient en multipliant le précédent par la même constante $1{,}03$.

\commentaireMarge{La relation $C_{n+1}=1{,}03C_n$ est précisément la définition d'une suite géométrique.}

La suite $(C_n)$ est donc géométrique de raison $q = 1{,}03$ et de premier terme $C_0 = 800$.

\medskip
\textbf{Question 2 — Expression de $C_n$ en fonction de $n$.}

Par la formule du terme général d'une suite géométrique :
\[
  C_n = C_0 \times q^n = 800 \times 1{,}03^n.
\]

\medskip
\textbf{Question 3 — Sens de variation.}

La raison est $q = 1{,}03 > 1$ et $C_n > 0$ pour tout $n \in \mathbb{N}$. Donc $C_{n+1} = 1{,}03 \times C_n > C_n$.

\commentaireMarge{$q > 1$ et termes positifs $\Rightarrow$ suite strictement croissante : le capital augmente chaque année.}

La suite $(C_n)$ est \textbf{strictement croissante} : le capital augmente chaque année. C'est cohérent avec le fait que les intérêts s'ajoutent au capital (intérêts composés).

\medskip
\textbf{Question 4 — Calcul de $C_5$.}

\[
  C_5 = 800 \times 1{,}03^5 \approx 800 \times 1{,}15927 \approx 927{,}42 \text{ \euro{}}.
\]

Après $5$ ans, Yasmine dispose d'environ $927{,}42$~\euro{}, soit un gain de $127{,}42$~\euro{} par rapport au capital initial.

\medskip
\textbf{Question 5 — Détermination du plus petit $n$ tel que $C_n > 1\,200$.}

On cherche le plus petit entier $n$ tel que :
\[
  800 \times 1{,}03^n > 1\,200 \iff 1{,}03^n > \frac{1\,200}{800} = 1{,}5.
\]

\commentaireMarge{On utilise les valeurs approchées données : $1{,}03^{13} \approx 1{,}469$ et $1{,}03^{14} \approx 1{,}513$.}

Par essais successifs, en utilisant les valeurs données :
\[
  1{,}03^{13} \approx 1{,}469 < 1{,}5 \implies C_{13} \approx 800 \times 1{,}469 = 1\,175{,}20 < 1\,200.
\]
\[
  1{,}03^{14} \approx 1{,}513 > 1{,}5 \implies C_{14} \approx 800 \times 1{,}513 = 1\,210{,}40 > 1\,200. \checkmark
\]

Ainsi $C_{13} < 1\,200 < C_{14}$. Le plus petit entier $n$ tel que $C_n > 1\,200$ est $\boxed{n = 14}$.

Yasmine doit attendre $14$ ans pour que son capital dépasse $1\,200$~\euro{}.

\baremeIndicatif{Q1 : 2 pts (modeliser, raisonner : relation de récurrence 1 pt + démonstration géométrique 1 pt) ; Q2 : 1 pt (calculer : terme général) ; Q3 : 2 pts (raisonner, communiquer : conclusion sur $q$ 1 pt + interprétation 1 pt) ; Q4 : 1 pt (calculer : $C_5$ arrondi) ; Q5 : 2 pts (raisonner, calculer : inéquation 1 pt + résolution par essais et conclusion 1 pt)}
\end{corrige}
